\documentclass[aspectratio=169,10pt,t]{beamer}
\usepackage{ac1}
\aulainfo{31}{Implementação 1}{Estrutura básica do programa, funções principais e dados}

\begin{document}
\capa

\begin{frame}{Objetivo da aula}
  \begin{itemize}
    \item Construir o \textbf{esqueleto} do programa: arquivos, funções principais e estruturas de dados
    \item Ter, ao final, uma versão que \alert{executa} do início ao fim (mesmo que incompleta)
    \item Feedback contínuo do professor durante a aula
  \end{itemize}
  \begin{block}{Princípio: desenvolvimento incremental}
    Escreva um pouco, \textbf{execute}, teste, corrija; só então avance. Nunca escreva 200 linhas antes de rodar pela primeira vez.
  \end{block}
\end{frame}

\begin{frame}{Roteiro de implementação}
  \begin{enumerate}
    \item Crie a \textbf{pasta e os arquivos} previstos na especificação
    \item Escreva os \textbf{esqueletos das funções} (nome, parâmetros, docstring, \texttt{pass})
    \item Defina a \textbf{estrutura de dados} central e crie dados de exemplo \emph{no próprio código}
    \item Implemente a \textbf{funcionalidade mais simples} de ponta a ponta (entrada \(\to\) cálculo \(\to\) saída)
    \item Escreva o \texttt{main()} com o menu/fluxo geral
    \item Teste, \textbf{integre} o trabalho dos integrantes e repita
  \end{enumerate}
\end{frame}

\begin{frame}[fragile]{Passo 2: esqueleto de funções}
\begin{lstlisting}[basicstyle=\ttfamily\scriptsize]
def cadastra_insumo(estoque, nome, unidade, preco, qtd, minimo):
    """Acrescenta um insumo ao estoque (dicionário). Devolve True se criou."""
    pass

def baixa_estoque(estoque, nome, qtd):
    """Retira qtd do estoque. Levanta ValueError se faltar material."""
    pass

def custo_total(estoque):
    """Soma preco * qtd de todos os insumos."""
    pass

def insumos_em_alerta(estoque):
    """Lista os nomes com qtd menor ou igual ao mínimo."""
    pass
\end{lstlisting}
  \small Nomes, parâmetros e \emph{docstrings} já definem o \textbf{contrato} de cada função. Os integrantes podem trabalhar em funções diferentes.
\end{frame}

\begin{frame}[fragile]{Passo 3 e 4: dados e primeira funcionalidade}
\begin{lstlisting}
def custo_total(estoque):
    """Soma preco * qtd de todos os insumos."""
    total = 0.0
    for dados in estoque.values():
        total += dados["preco"] * dados["qtd"]
    return total

def main():
    estoque = {
        "cimento": {"unidade": "saco", "preco": 38.5, "qtd": 120, "minimo": 30},
        "areia":   {"unidade": "m3",   "preco": 95.0, "qtd": 15,  "minimo": 10},
    }
    print(f"Custo total em estoque: R$ {custo_total(estoque):.2f}")

main()
\end{lstlisting}
\end{frame}

\begin{frame}[fragile]{Passo 5: o menu principal}
\begin{lstlisting}[basicstyle=\ttfamily\scriptsize]
def mostra_menu():
    print("\n1 - Cadastrar insumo\n2 - Baixar estoque")
    print("3 - Custo total\n4 - Alertas\n0 - Sair")

def main():
    estoque = {}
    while True:
        mostra_menu()
        opcao = input("Opção: ").strip()
        if opcao == "1":
            ...      # le dados e chama cadastra_insumo
        elif opcao == "3":
            print(f"Custo: R$ {custo_total(estoque):.2f}")
        elif opcao == "0":
            break
        else:
            print("Opção inválida.")
\end{lstlisting}
  \small Cada opção do menu chama uma função: \texttt{main} só \emph{coordena}.
\end{frame}

\begin{frame}{Dividindo o trabalho em equipe}
  \begin{columns}[T]
    \column{0.5\textwidth}
    \begin{block}{Boas práticas}
      \begin{itemize}\small
        \item Cada integrante assume \textbf{funções ou módulos} distintos
        \item Combinem o \textbf{contrato} (parâmetros e retorno) antes de escrever
        \item Evitem editar o mesmo arquivo ao mesmo tempo
        \item Integrem \emph{no mesmo dia}, não na véspera
      \end{itemize}
    \end{block}
    \column{0.5\textwidth}
    \begin{alertblock}{Armadilhas}
      \begin{itemize}\small
        \item Cada um com sua versão ``do arquivo''
        \item Variáveis globais compartilhadas entre módulos
        \item Funções que leem do teclado \emph{e} calculam (difíceis de testar)
        \item Nomes diferentes para a mesma coisa
      \end{itemize}
    \end{alertblock}
  \end{columns}
\end{frame}

\begin{frame}{Estruturas de dados: escolha consciente}
  \centering\small
  \begin{tabular}{p{4.6cm}p{6.6cm}}
    \toprule
    \textbf{Preciso de\ldots} & \textbf{Use}\\
    \midrule
    uma coleção ordenada de itens iguais & lista\\
    buscar por nome/código & dicionário\\
    registro fixo (coordenada, data) & tupla\\
    muitas operações numéricas & \texttt{np.array}\\
    tabela com colunas nomeadas & lista de dicionários ou \texttt{DataFrame}\\
    grade/tabuleiro & lista de listas (matriz)\\
    \bottomrule
  \end{tabular}
  \vspace{1ex}

  \footnotesize ``Uma estrutura de dados mal escolhida pode arruinar o melhor algoritmo.'' (Cap.~9.1)
\end{frame}

\begin{frame}{Atividade: ao final desta aula}
  Cada grupo deve mostrar ao professor:
  \begin{itemize}
    \item A \textbf{estrutura de pastas/arquivos} criada
    \item Os \textbf{esqueletos das funções} com docstrings
    \item Pelo menos \textbf{uma funcionalidade} completa funcionando (entrada \(\to\) saída)
    \item O \texttt{main()} com o fluxo geral (menu ou sequência de passos)
  \end{itemize}
  \begin{exampleblock}{Feedback do professor}
    Verificará: funções coesas? Nomes claros? Dados bem escolhidos? O programa \emph{executa}?
  \end{exampleblock}
\end{frame}

\begin{frame}{Checklist de implementação 1}
  \small
  \begin{tabular}{@{}cl@{}}
    $\square$ & O programa roda sem erro do início ao fim?\\
    $\square$ & Cada função tem nome, parâmetros e docstring?\\
    $\square$ & Há uma função \texttt{main()} e o programa está dividido em módulos/funções?\\
    $\square$ & A estrutura de dados central está definida?\\
    $\square$ & Uma funcionalidade essencial está completa?\\
    $\square$ & O que falta está listado (para a Aula 32)?\\
  \end{tabular}
\end{frame}

\end{document}
